\section{Thirteen integers without the 2026 Sun reduction}
There is an explicit polynomial in thirteen unrestricted integer unknowns of exact degree $10\,461\,744$, evaluated by 519 binary arithmetic operations, whose absence of integer zeros is equivalent in ZFC to $\operatorname{Con}(\mathrm{ZFC})$. This is a degree/cost tradeoff with the preceding $11\,543\,632$-degree, 518-operation form, not a strict operation-count improvement. The unchanged finalizer uses the Jones ratio theorem, not Sun's 2026 reduction. The new source uses the split108 placement and a lower off-target positivity pad. A necessary new parity adjustment makes the square threshold cubic in $\epsilon$ rather than linear.

\subsection{Degree-separated kernels and the complete polynomial}
Use the original package-31 system of 22 quadratic residuals
$r_j(p)$ in 31 strictly positive coordinates. The exact data are
\texttt{split/parent\_source.json}; the coherent fixed port tuple is
$(\beta,K_0,M_C,M_F)=(\beta,K,M_C,M_{F,\mathrm{source}})$ from
\texttt{FIXED\_BINDING.json}. All four compressed little-endian payloads are
included and hash-authenticated. The inherited compiler theorem identifies
its positive common zeros with finite syntactic ZFC refutations. This is the
original tuple, not a mixture with a later packet's smaller coefficients.

Perform just two elementary integer row operations:
\[
 \varrho_{14}=-(r_{14}+r_2),\qquad
 \varrho_{15}=-(r_{15}+r_3),\qquad \varrho_j=-r_j\quad(j\ne14,15).
\]
First negate all rows and then undo the two shears, so the common zero set is
unchanged. Write $\varrho_j=\sum_\alpha r_{j,\alpha}p^\alpha$ and
\[
 c_{j,\alpha}=\begin{cases}
 2r_{j,\alpha},&|\alpha|=0,1,\\
 (\prod_i\alpha_i!)r_{j,\alpha},&|\alpha|=2.
 \end{cases}
\]
Use the 31 positive weights $w_i$ and the 22 targets $q_j$ in
\texttt{split/split\_code.json}, and set
\[
 D_d(T)=\sum_{j,\,|\alpha|=d}c_{j,\alpha}T^{q_j-\alpha\cdot w}
 \quad(d=0,1,2).
\]
The supplied exact certificate establishes
\[
 W=\max_iw_i=108,\quad Q_*=\max_jq_j=1939,\quad
 (\deg D_0,\deg D_1,\deg D_2)=(1939,1842,1738).
\]
The leading coefficients of $(D_0,D_1,D_2)$ are $(2,2,\beta)$; every kernel has internally
distinct exponents. There are 89 source records. Checking all 22,544
same-degree source/expansion interactions proves
\[
 [T^{q_j}]\bigl(D_2(T)U(T)^2+D_1(T)U(T)+D_0(T)\bigr)
 =2\varrho_j(p),\qquad U(T)=\sum_i p_iT^{w_i}.
\]
The degrees are kept in separate products. Thus a linear label may equal a
quadratic label without causing cross-degree contamination. No homogeneous
dummy coordinate or full Sidon condition is needed. Let
\[
 K_*=\max_d(\deg D_d+dW)=1954,\qquad E_*=K_*+W+2=2064.
\]
The fixed coefficient bound, checked from the actual original port bit
lengths, is $|c_{j,\alpha}|<A_*=2^{2134927}$.

The unknowns are exactly
\[
 (\epsilon,g,h,s,w,\phi,i,j,z_1,z_2,z_3,r,v)\in\mathbb Z^{13}.
\]
All following expressions are polynomial abbreviations. Fix
\[
 \kappa=2134940,\qquad a=2^\kappa=8192A_*,\qquad
 \eta=2^{18146981}=2^{17\kappa/2-9},
\]
and define
\begin{align*}
 b&=2\epsilon^2,& B&=ab^2,\\
 \ell&=\sum_iB^{w_i},& T_*&=\sum_jB^{q_j},\\
 U_c&=\ell+g,&\mathcal G&=D_2(B)U_c^2+D_1(B)U_c+D_0(B),\\
 \mathcal G_+&=\mathcal G+B^{1946},&N_0&=B^{109},\\
 L&=B^{2064},&\widehat N&=\eta\epsilon^3B^{1089},\\
 c&=bB^{108},&G_c&=32c,\\
 M_c&=N_0-1-(b-1)\ell,&J_c&=2\mathcal G_++T_*,\\
 S_c&=g+N_0J_c,&T_c&=M_c+N_0(B-2)T_*,\\
 R&=L^2+(S_c+T_c)L+(L-1-S_c).
\end{align*}
The extra monomial $B^{1946}$ lies above every target. It repairs positivity
on arbitrary signed $g$ without changing an extracted residual or adding a
new unknown. Every exponent in these formulas is fixed and nonnegative.

The circuit's register named $N$ means the signed polynomial $\widehat N$. In the proof only, let $N=|\widehat N|$. Every occurrence of this parameter in the equation is through $N^2=\widehat N^2$, so no absolute-value operation occurs in the polynomial.

Use the Jones registers with the exact CRT reparameterization of $H$ below, and $\omega=w^2$:
\begin{align*}
 U&=N^2\omega,&Y&=N^2s,\\
 M&=RY,&P&=2M^2U,\\
 K&=R+1+h(P-1),&A&=M(U+1),\\
 B_1&=2R+1,&W_1&=b\omega,\\
 C&=B_1+W_1+\phi,&J&=4A-5,&
 Q&=3W_1C-2(W_1^2-1),\\
 D&=(A^2-1)C^2+1,&E&=iJDC^2,&F&=(A^2-1)E^2+1,\\
 G&=A+F(CD+1-A),&H&=C+F(B_1+jC),&I&=(G^2-1)H^2+1,\\
 a_0&=(P^2-1)K^2+1,&a_1&=DFI,\\
 u&=J,&t&=Q.
\end{align*}
Set
\[
 \delta=C-KY,\qquad \Delta=g(c-g)K^2-c^2\delta^2,
 \qquad (d_1,d_2,d_3)=(3s-1,3\phi+2,3\Delta-1).
\]
The new coprime-spectator stack is
\[
 C_0=a_0,\qquad B_j^*=1+4d_jC_{j-1}z_j^2,\qquad
 C_j=C_{j-1}B_j^*\quad(j=1,2),
\]
\[
 B_3^*=1+4d_3C_2a_1z_3^2,\qquad C_*=C_3=C_2B_3^*.
\]
Thus the last guard contains the second square candidate $a_1$ as a
spectator. This is the only change to the Pell sign stack.
Finally define
\begin{align*}
 Z&=\epsilon wi z_1z_2z_3a_0,&\Sigma&=Zu^5,&\rho&=t+ur,\\
 \Theta&=(\rho^3-u^2\rho-u^3)(\rho^2+23u^2),&
 X&=\Sigma v-\Theta,\\
 A_0&=\Sigma^2C_*,&A_1&=\Sigma^2a_1.
\end{align*}
The requested thirteen-integer equation is the following single polynomial:
\begin{equation}\label{eq:jones-seeded}
 \boxed{\mathcal F_{13}=(X^2-A_0-A_1)^2-4A_0A_1=0.}
\end{equation}
The literal division-free DAG is \texttt{generated/signed\_coprime\_compact\_split\_crt.json}.

\subsection{Signed-input bounds and exact square threshold}
Assume only $\epsilon\ne0$, hence $b\ge2$ and $B=8192A_*b^2$. Since all lower coefficients are bounded by $A_*$ and the leaders are $(2,2,\beta)$, geometric summation gives
\[
D_0(B)>0,\qquad 0<D_1(B)<3B^{1842},\qquad D_2(B)>(\beta-1)B^{1738}.
\]
Here $\beta=2^{1064921}-1>10$. Completing half the quadratic square, for every real $U_c$,
\[
\mathcal G_+\ge\tfrac12D_2(B)U_c^2+B^{1946}-\frac{D_1(B)^2}{2D_2(B)}
>\tfrac12D_2(B)U_c^2+\tfrac12B^{1946},
\]
because the subtracted term is less than $9B^{1946}/(2(\beta-1))<B^{1946}/2$.
Thus for every integer $g$, writing $U_c=\ell+g$,
\[
S_c>N_0D_2(B)U_c^2+U_c+N_0B^{1946}-\ell>0.
\]
The first two terms have nonnegative sum for integer $U_c$, and $\ell<N_0$. The mask gives $T_c>0$ as well. Therefore $R>L^2$ before any sign guard or source decoding.

The new threshold is not the old linear-$\epsilon$ expression. It satisfies the exact identity
\[
\widehat N^2=2^{-21}a^{17}b^3B^{2178}<B^{2195}<B^{4128}=L^2.
\]
Also $N=|\widehat N|=\eta|\epsilon|^3B^{1089}\ge bB^{108}=c$, since $\eta|\epsilon|B^{981}\ge2$, and $N\ge8$. Thus the inherited ratio hypotheses $R\ge N\ge8$ and $b,c\le N$ hold before signs of the other free variables are recovered.

For $0<g<c$, $U_c<2bB^{108}$ and the old coefficient-majorant argument gives
\[
|\mathcal G|<12A_*b^2B^{1954},\qquad
2|\mathcal G|<B^{1955}/256,\qquad T_*<B^{1955}/256.
\]
The extra pad $B^{1946}$ is smaller than $B^{1954}$, so $0<J_c<B^{1955}/64$. Consequently $S_c<L/32$ and $T_c<L/64$, in particular $S_c+T_c<L$. The pad is above every target, since $1946>1939$, and below the convolution height. It changes no tested residual. All unnormalized coefficients of $\mathcal G_+$ remain below $B/8$ in absolute value under the inherited $994A_*b^2$ bound.

Exactly the sparse-borrow argument gives
\[
S_c\mathbin{\&}T_c=0\quad\Longleftrightarrow\quad
 g=\sum_i n_iB^{w_i},\quad0\le n_i<b,\quad r_j(n+\mathbf1)=0\ (1\le j\le22).
\]
Indeed the unnormalized target digit is $4\varrho_j(n+\mathbf1)+1$ and the incoming borrow is $0$ or $-1$; it avoids the mask $B-2$ exactly when that residual vanishes. Non-target collisions cause no problem. The source all-one vector is not a zero, so the forward encoded $g$ is positive. Cofinal choices $b=2^{1+2r}$ accommodate every positive source tuple.

After the Jones theorem recovers $b=2^k$, the identity $b=2\epsilon^2$ forces odd $k$. Put $B=2^\nu$, $\nu=\kappa+2k$. The actual mask population is $131\nu-31k-22$, and the tight source range yields
\[
s_2(R)=2195\nu-31k-21-\tau_2(S_c,T_c).
\]
The cubic-$\epsilon$ threshold is calibrated exactly:
\[
2195\nu-31k-21
=2(17\kappa/2-9+1089\nu)+3(k-1)
=\log_2\widehat N^2.
\]
The exponent $3(k-1)$ comes from $\epsilon^6$, not $\epsilon^2$; retaining the previous linear-$\epsilon$ threshold would be incorrect for this source parity. Legendre's identity therefore proves $N^2\mid\binom{2R}{R}$ exactly when the masks are disjoint. The following retained finalizer proof now applies without any other source change.

\subsection{An exact polynomial CRT reparameterization}
The inherited ratio interface uses $H_{\rm old}=B_1+j_{\rm old}C$ and
$JF\mid QF+J(H_{\rm old}-C)$. Since $E=iJDC^2$,
\[
 F\equiv1\pmod J,\qquad F\equiv1\pmod C,
\]
and hence $\gcd(J,F)=1$. The divisibility is equivalent to
$J\mid Q$ and $F\mid H_{\rm old}-C$. Our changed register
\[
 H=C+F(B_1+jC)
\]
automatically satisfies the latter. It also has the original form, with
\[
 j_{\rm old}=1+(A^2-1)i^2J^2D^2C^3B_1+Fj,
\]
which is a polynomial integer, since $F-1=(A^2-1)i^2J^2D^2C^4$.
This identity holds for arbitrary signed inputs, including $C=0$.

Conversely, at the canonical positive ratio witnesses, $0<B_1<C$ and
$j_{\rm old}\ge1$. Therefore $z=(H_{\rm old}-C)/F>0$ and
$z\equiv B_1\pmod C$, so $z=B_1+jC$ with $j\ge0$. Thus every
canonical witness is represented. No extra unknown or cofinal restriction is
introduced. The merger now has $u=J$ rather than $JF$. This CRT step is
inherited from the separately identified sparse-Pell packet, not claimed as
an original theorem here.

\subsection{The interval guard and the Jones interface}
If $\Delta>0$ and $b\ge2$, then $K\ne0$ and
$0<g<c$. The identity $4g(c-g)=c^2-(c-2g)^2\le c^2$ implies
$4\delta^2<K^2$. Along with $s\ge1$, $\phi\ge0$, and $w\ne0$, this gives
$C,Y>0$ and $K>0$. The usual square-gap argument forces $h\ge1$:
if $h\le0$, then $1\le K\le R+1<P$; $K=1$ would force $P-1\mid R$,
which is impossible, while for $2\le K<P$,
\[
 (PK-1)^2<(P^2-1)K^2+1<(PK)^2.
\]
Thus $a_0$ cannot be a square in that branch.

For precision, instantiate the free ceiling in
\texttt{Jones1982.Nine.ratio\_sound} by
$G_{\mathrm{aux}}=8gK^2+1$, rather than assuming its old margin condition.
The positive integer $K^2-4\delta^2$ then gives
\[
 (K^2-4\delta^2)G_{\mathrm{aux}}-8gK^2>0.
\]
This ceiling is a proof-side abbreviation; the ratio registers do not
otherwise depend on it.

Replace $i$ by $|i|$, which changes no condition. To handle signed $j$,
$u\mid t$ now means $J\mid Q$, while $F\mid H-C$ holds identically. Also
$F\equiv1\pmod D$, $G\equiv1\pmod D$, and $G\equiv A\pmod F$ show
that the positive integers $D,F,I$ are pairwise coprime.
Hence $a_1=DFI$ square makes each a square. Put $d=G^2-1>0$ and
$x=\sqrt I$. The norm-one unit $x+H\sqrt d$ is positive even if $H<0$.
A positive power of $G+\sqrt d$ is the identity modulo $CF$; multiplying by
sufficiently large such powers gives a Pell solution $H'>B_1$ with
$H'\equiv H\pmod{CF}$. Then $j'=(H'-B_1)/C>0$, all square conditions
remain true, and $J\mid Q$ remains true. The literal positive Jones interface
therefore applies and gives the power-of-two base and central-binomial
condition required above.

For completeness, take any common source zero, choose
$b=2^{1+2r}$ large enough, and set $\epsilon=2^r$. The direct source code
has $32g<G_c$, so the literal ratio completeness theorem supplies its
canonical witnesses. In \texttt{Ratio.lean}, \texttt{exists\_B25core} sets
\[
 \omega=2^{2R+1-k},\qquad U=(N^2/b)2^{2R+1}.
\]
Since $k=1+2r$ is odd and $k<2R+1$, this $\omega$ is a positive square.
No claim that the canonical $s$ is a square is needed.

The same canonical construction has
$\rho_0=(U+1)^{2R}/U^R=Y+f$ with
$0\le f<2^{2R}/U=b/(2N^2)$ and $Y\ge1$.
The source's \texttt{ratio\_core} gives
\[
 |C/K-\rho_0|\le\frac{2(Y+f)}{Y(U+1)}<\frac4U,
 \qquad |C/K-Y|<\frac b{N^2}.
\]
The last inequality uses $R\ge8$. Since $c,b\le N$ and
$g(c-g)\ge1$, it follows that
\[
 \frac{c^2(C-KY)^2}{K^2}<\frac{c^2b^2}{N^4}\le1\le g(c-g),
\]
so the very same canonical witnesses, translated by the exact CRT formula, satisfy $\Delta>0$.
This is a cofinal square-base proof, not an unsupported restriction of an
existing slack variable to squares.

\subsection{Raw two-root recovery, including the zero spectator branch}
For every integer assignment, $u=J=4A-5\ne0$. The polynomial
\[
 \Pi(x)=(x^3-x-1)(x^2+23)
\]
has no rational root and has a root modulo every positive integer.
For an odd prime other than $23$, either the cubic has a simple root or it
is irreducible, whose Frobenius three-cycle makes its nonzero discriminant
$-23$ a square. The quadratic then has a simple root. At $23$ the cubic has
the simple root $3$; at $2$, $-23\equiv1\pmod8$ admits compatible square
roots. Hensel lifting and CRT finish the assertion. Thus
$\Theta=u^5\Pi(t/u+r)\ne0$ for every assignment, and
\[
 \exists r\quad Zu^5\mid\Theta
 \quad\Longleftrightarrow\quad Z\ne0\ \hbox{and}\ u\mid t.
\]
Indeed a root modulo $|Z|$ proves the reverse implication once $u\mid t$.
If $v_p(u)=\alpha>\beta=v_p(t)$, the two factors of $\Theta$ have
valuations $3\beta$ and $2\beta$, contrary to $u^5\mid\Theta$.
If $Z=0$, divisibility would require $\Theta=0$, already impossible.

At a zero of the equation, $\Sigma=0$ is impossible because then the
polynomial is $\Theta^4\ne0$. In particular
$\epsilon,w,i,z_1,z_2,z_3,a_0$ are nonzero. Every $d_j\equiv2\pmod3$
is nonzero, and every $B_j^*\equiv1\pmod4$ is nonzero. Thus $C_2,C_*$
are nonzero.

\paragraph{No coincident radicands.}
One has $C_*\ne a_1$ without any positivity assumption. If $a_1=0$ this
follows from $C_*\ne0$. Otherwise $B_3^*\ne1$, and it cannot be $-1$
because it is $1$ modulo $4$. But
\[
 \gcd(B_3^*,a_1)=1,
\]
while $C_*=a_1$ would make $B_3^*$ divide $a_1$. This would force
$B_3^*=\pm1$, a contradiction.

Put $q=X/\Sigma$. The raw norm equation implies
\[
 (q^2+C_*-a_1)^2=4C_*q^2.
\]
If $q=0$, the norm would give $C_*=a_1$, just excluded. Therefore
\[
 r_0=\frac{q^2+C_*-a_1}{2q},\qquad r_1=q-r_0
\]
are rational, with $r_0^2=C_*$ and $r_1^2=a_1$. Rational square roots of
integers are integers. Hence $C_*,a_1$ are integer squares,
$q\in\mathbb Z$, and $\Sigma\mid\Theta$. In particular $J\mid Q$.
No complex-root separation assumption is used.

\paragraph{Recover the first two guards before ruling out $a_1=0$.}
The factors $a_0,B_1^*,B_2^*,B_3^*$ are pairwise coprime, because each
new factor is $1$ modulo its entire preceding product. Their nonzero product
$C_*$ is a square. Thus each absolute factor is a square. Since
$B_j^*\equiv1\pmod4$, every $B_j^*$ is positive: a negative square cannot
be $1$ modulo $4$. The product sign then forces $a_0>0$, also a square.
All prefixes are now positive, so $d_1,d_2>0$, yielding $s\ge1$ and
$\phi\ge0$.

At this point the last factor might still equal $1$ if $a_1=0$; this branch
must not be silently discarded. The unconditional source estimate above
has already proved $R>L^2>0$ for \emph{every} signed $g$. Since
$\epsilon,w\ne0$ and $s\ge1$, we have $b\ge2$, $N\ge8$, $U,Y>0$,
$A\ge2$, and, using $\phi\ge0$, $C>0$. Then
\[
 D\ge A^2>0,\qquad F\ge1,\qquad
 CD+1-A>0,\qquad G\ge A>1,\qquad I\ge1.
\]
Consequently $a_1=DFI>0$. Now $B_3^*\ne1$, and its positivity with
$C_2a_1z_3^2>0$ forces $d_3>0$, hence $\Delta>0$. This closes the
zero-spectator branch without adding $a_1$ to the nonzero multiplier.
The interval and ratio arguments above therefore apply.

\paragraph{Completeness.}
Start from the canonical positive ratio witnesses, translated by the CRT
parameterization. They have $a_0,a_1$ positive squares, $J\mid Q$, and all
three strict sign conditions. For $j=1,2$, $d_jC_{j-1}$ is positive
nonsquare because $d_j\equiv2\pmod3$ and $C_{j-1}$ is a square.
A nontrivial Pell solution $x^2-d_jC_{j-1}y^2=1$ can be chosen with $y$
even, so $z_j=y/2\ne0$ makes $B_j^*=x^2$. For the last guard,
$d_3C_2a_1$ is likewise positive nonsquare, since $C_2a_1$ is a positive
square. The same construction gives $B_3^*$ square and $C_*$ square.
The intersective lemma chooses $r$ with $\Sigma\mid\Theta$. Finally choose
\[
 v=\Theta/\Sigma+\sqrt{C_*}+\sqrt{a_1}\in\mathbb Z.
\]
Then the raw norm vanishes. Both directions are proved. By the inherited
compiler theorem, integer-zero existence is equivalent to a finite ZFC
refutation, and ZFC proves its negation equivalent to
$\operatorname{Con}(\mathrm{ZFC})$.

\subsection{Exact degree and validation}
The emitted DAG is \texttt{generated/signed\_coprime\_compact\_split\_crt.json}. It uses thirteen unrestricted variables and 519 live gates: 327 multiplications, 169 additions, and 23 subtractions. Fixed leaf construction is excluded from this arithmetic metric. The supplied four coefficient integers are unchanged.

The independent verifier propagates an upper degree and the leading coefficient on the full diagonal using the actual fixed coefficients. The leading residues at degree $10461744$ are $20$, $442$, and $336577$ modulo $101$, $1009$, and $1000003$, respectively. Any one nonzero residue certifies that this upper degree is attained. It also replays all named registers against independent formulas on signed assignments, checks the source certificate, and tests the guard and CRT identities. The replay is finite supplementary evidence, not a re-proof of the imported unbounded compiler or Jones theorem. The full retained prior proof and provenance are included under \texttt{upstream/}; the active definitions and cubic threshold above supersede the prior source registers.
